\subsection{A small cyclic part, and what remains missing on general graphs}
\label{sec:threshold-positive-scope}

\begin{corollary}[Bounded feedback vertex sets]
\label{cor:threshold-feedback-sets}
Fix \(r<\infty\). Suppose a known set \(S\) of at most \(r\) vertices can be deleted
from an Ising interaction graph to leave a forest. At every fixed
\(\beta J_{\mathrm{edge}}<\infty\), its Gibbs laws have full AA with the query
bound in Theorem~\ref{thm:threshold-ising-forests} increased by \(r\).
The fields, degrees, and couplings incident to \(S\) may vary
arbitrarily; the unpinned forest edges must satisfy the stated
\(J_{\mathrm{edge}}\) bound.
\end{corollary}

\begin{proof}
Query all of \(S\). Conditional on its revealed assignment, the
remaining Hamiltonian is an Ising forest Hamiltonian with additional
fields. Fixing the \(S\)-coordinates in the quantum oracle leaves a
quantum algorithm with at most \(T\) queries to the remaining input.
Choose the forest simulator for each positive-probability assignment
of \(S\). The field-uniform query bound is common to all branches,
and averaging their conditional errors gives the same mean error.
\end{proof}

\textbf{Structural implication.} A large tree, long correlation length,
or bounded-size cyclic separator is not by itself a source of a
finite-temperature full-AA obstruction. On a tree, the important fact
is that the randomness generating dependence has an edge-local
posterior, so it can be recovered lazily from the observed input. This
remains true even when a root is correlated with many descendants and
local flip-influence statements are inappropriate proxies.

\textbf{Research direction, not a theorem for all high-temperature
Gibbs laws.} General graphs suggest replacing the single ancestor
path by a random dependency cluster. A sufficient route would be:
\begin{enumerate}
\item a source representation by independent bounded-size blocks;
\item coherent evaluation after truncation using
\(\operatorname{poly}(T,1/e)\) source queries;
\item conditional, rather than merely marginal, disagreement at most
\(e^2/(64T^2)\);
\item a posterior source oracle that answers the polynomially many
adaptive requests with polynomially many input queries.
\end{enumerate}
Under those four conditions the transfer proposition proves full AA.
A physical uniqueness or percolation threshold is relevant when it
controls these quantitative properties. It is not itself the missing
query-theoretic implication.

There is substantial precedent for the forward part:
\href{https://arxiv.org/abs/1803.10578}{Spinka, \emph{Finitary codings
for spatial mixing Markov random fields}} proves that exponential
strong spatial mixing gives an i.i.d. coding with exponential coding
radius tails. This supplies a reason to investigate such encodings;
it does not assert the posterior-access or pointwise conditional
tail bounds used here. Nor does an exponential radius tail by itself
give a dimension-free query bound on graphs with rapidly growing
balls without controlling the constants and volume growth. The
forest proof establishes all of these query-relevant steps directly.

\textbf{Status and provenance.} The local posterior criterion,
atom-independent block theorem, forest theorem, and feedback-set
corollary are deductions recorded in this investigation. The
coordinatewise quantum hybrid bound and the uniform-AA simulation
input were already used in the shared exchangeable-simulation notes;
conditional gadget sampling is also used in
\href{https://arxiv.org/abs/2609.35327}{Girish--Mittal--Nehoran--Raz}.
No novelty or research-priority claim is made. The criterion on general
dependency clusters is a proposed proof program, not a classification
of arbitrary pair interactions or noncommuting Hamiltonians.
